\documentclass[11pt]{article}
\usepackage[T1]{fontenc}
\usepackage{lmodern}
\usepackage[margin=1in]{geometry}
\usepackage{amsmath,amssymb,amsthm,mathtools,booktabs}
\usepackage[hidelinks]{hyperref}
\newtheorem{theorem}{Theorem}
\newtheorem{lemma}[theorem]{Lemma}
\newtheorem{corollary}[theorem]{Corollary}
\theoremstyle{remark}
\newtheorem{remark}[theorem]{Remark}
\newcommand{\HH}{\mathfrak H}
\newcommand{\SL}{\mathrm{SL}}
\newcommand{\PSL}{\mathrm{PSL}}
\newcommand{\eps}{\varepsilon}
\newcommand{\bb}{\beta}
\title{A twist-free converse theorem at level $36$}
\author{}
\date{27 October 2026}

\begin{document}
\maketitle
\begin{abstract}
We prove the holomorphic Conrey--Farmer converse theorem at level $36$,
for every positive even weight and either functional-equation sign.
Only the local Euler conditions at $2,3,5$ are needed.  For the positive
Fricke eigenvalue, the conditions at $2,3$ alone suffice.  The proof uses
the Fourier support forced by the ramified Euler factors to construct a
three-dimensional module for the full modular group.  A single $m=5$
relation eliminates its possible noncongruence part.  No twisted
functional equations or additional coefficient estimates are assumed.
\end{abstract}

\section{Statement and conventions}

Conrey--Farmer \cite{CF} establish their holomorphic converse theorem for
$5\le N\le12$, $14\le N\le17$, and $N=23$.  Conrey--Farmer--Odgers--Snaith
\cite{CFOS} supply $N=13$, and Bedert--Cooper--Oliver--Zhang
\cite{BCOZ} supply $N=18,20,24$.  The theorem below adds $N=36$ to this
list.

\begin{theorem}\label{thm:main}
Let $k\ge2$ be even.  Let $a_1=1$, and suppose
\[
 F(s)=\sum_{n\ge1}\frac{a_n}{n^s}
\]
converges in a right half-plane.  Assume
\begin{enumerate}
\item $\Lambda(s)=(6/(2\pi))^s\Gamma(s)F(s)$ extends to an entire function,
bounded in every closed vertical strip, and
\[
 \Lambda(s)=\omega\Lambda(k-s),\qquad \omega\in\{1,-1\};
\]
\item $a_n=0$ whenever $(n,6)>1$;
\item for all $n\ge1$,
\[
 a_{5n}=a_5a_n-5^{k-1}a_{n/5},
\]
where $a_{n/5}=0$ if $5\nmid n$.
\end{enumerate}
Then
\[
 f(z)=\sum_{n\ge1}a_n e^{2\pi inz}\in S_k(\Gamma_0(36)).
\]
If $(-1)^{k/2}\omega=1$, hypothesis~(3) is unnecessary.
\end{theorem}

The local hypotheses in the theorem follow from the partial Euler product
\[
 F(s)=\frac{1}{1-a_5 5^{-s}+5^{k-1-2s}}
       \sum_{(n,30)=1}\frac{a_n}{n^s}.
\]
In particular, they follow from the full Conrey--Farmer Euler product at
level $36$: its factors at $2$ and $3$ are $1$.  No condition is imposed
here at any other prime.

For $M=\begin{pmatrix}a&b\\c&d\end{pmatrix}\in\mathrm{GL}_2^+(\mathbb R)$,
use the normalized weight-$k$ action
\[
 (v|_k M)(z)=(\det M)^{k/2}(cz+d)^{-k}
 v\left(\frac{az+b}{cz+d}\right).
\]
We suppress the subscript $k$.  Throughout $k$ is positive and even, so
nonzero scalar matrices act trivially and $-I$ acts trivially.  Set
\[
 S=\begin{pmatrix}0&-1\\1&0\end{pmatrix},\quad
 T=\begin{pmatrix}1&1\\0&1\end{pmatrix},\quad
 \bb(t)=\begin{pmatrix}1&t\\0&1\end{pmatrix},\quad
 H_N=\begin{pmatrix}0&-1\\N&0\end{pmatrix}.
\]

Convergence in a right half-plane implies $a_n=O(n^C)$ for some $C\ge0$:
choose a real point of convergence, use that its summands are bounded,
and increase the exponent if necessary.  Thus the Fourier series defines
a holomorphic function on $\HH$.  The usual Mellin-inversion argument,
as in \cite[Section~1]{CF}, gives
\begin{equation}\label{eq:fricke}
 f|H_{36}=\eps f,\qquad \eps=(-1)^{k/2}\omega\in\{1,-1\}.
\end{equation}
We distinguish the analytic root number $\omega$ from the Fricke
eigenvalue $\eps$ to avoid a weight-dependent sign ambiguity.

\section{Two elementary lemmas}

\begin{lemma}[Positive-weight elliptic vanishing]\label{lem:elliptic}
Let $v$ be holomorphic on $\HH$, and let $k>0$ be even.  If
$v|E=v$ for an infinite-order elliptic element $E\in\SL_2(\mathbb R)$,
then $v=0$.
\end{lemma}
\begin{proof}
Let $a\in\HH$ be the fixed point.  Use the disk coordinate
$w=(z-a)/(z-\bar a)$ and write $m=k/2$.  The holomorphic differential
$v(z)(dz)^m$ has the form $V(w)(dw)^m$, with $V$ holomorphic on the disk.
The action of $E$ is $w\mapsto\zeta w$, where $|\zeta|=1$ and $\zeta$
has infinite order.  Invariance says
$V(w)=\zeta^m V(\zeta w)$.  If $V(w)=\sum_{j\ge0}c_jw^j$, then
$(1-\zeta^{j+m})c_j=0$.  Since $j+m>0$, every coefficient vanishes.
\end{proof}

\begin{lemma}[A congruence-kernel criterion]\label{lem:kernel}
Let $V$ be a complex right $\PSL_2(\mathbb Z)$-module such that $T^3$
acts as $-\mathrm{Id}_V$.  Then $\Gamma(6)$ acts trivially on $V$.
\end{lemma}
\begin{proof}
Let $\chi_2$ be the sign of the permutation induced by reduction modulo
$2$ on $\mathbb P^1(\mathbb F_2)$.  Thus
$\chi_2(S)=\chi_2(T)=-1$, and $\chi_2$ is trivial on $\Gamma(2)$.
Twist the action by $\chi_2$.  In the twisted action, the operators
corresponding to $S,T$ satisfy
\[
 S^2=1,\qquad (ST)^3=1,\qquad T^3=1.
\]
The quotient of $\PSL_2(\mathbb Z)$ obtained by imposing $T^3=1$ is the
triangle group $(2,3,3)$, of order $12$, and reduction modulo $3$
identifies it with $\PSL_2(\mathbb F_3)\simeq A_4$.  For completeness,
this is the rotational symmetry group of a tetrahedron: its stabilizers
at an edge midpoint, a vertex, and a face center have orders $2,3,3$.
Equivalently, its standard triangle presentation gives the above
relations and order $12$.  Reduction modulo $3$ is surjective onto a
group of that same order, so its kernel is exactly the normal closure
of $T^3$ in the projective modular group.
Consequently the twisted action is trivial on $\Gamma(3)$.
On $\Gamma(6)\subseteq\Gamma(2)\cap\Gamma(3)$ the twist itself is trivial,
proving the assertion.
\end{proof}

\section{The three-dimensional module}

Rescale the function by setting
\[
 h(z)=f(z/6).
\]
Equation~\eqref{eq:fricke} gives $h|S=\eps h$.  Since the only possible
Fourier indices of $f$ are congruent to $1$ or $5$ modulo $6$, every
frequency of $h$ has translation eigenvalue a primitive sixth root.
Therefore
\begin{equation}\label{eq:phi6}
 h|(T^2-T+1)=0,
 \qquad h|T^{-1}=h-h|T,
 \qquad h|T^3=-h.
\end{equation}
Define
\[
 g=h|T,\qquad v=g|S.
\]
Using $STS=T^{-1}ST^{-1}$, we have
\[
 \eps v=h|STS
     =h|T^{-1}ST^{-1}
     =(\eps h-v)|T^{-1}.
\]
Right-multiplication by $T$ yields $v|T=h-\eps v$.  Hence
\begin{equation}\label{eq:table}
\begin{array}{c|cc}
 & |S & |T\\ \hline
 h & \eps h & g\\
 g & v & g-h\\
 v & g & h-\eps v
\end{array}
\end{equation}
In particular, $V=\operatorname{span}_{\mathbb C}\{h,g,v\}$ is stable
under $S$ and $T$, and therefore under the full modular group.
Linear independence of these three functions is not assumed or needed.

The formal matrices for \eqref{eq:table}, with functions as a column,
are
\begin{equation}\label{eq:matrices}
 A_\eps=\begin{pmatrix}\eps&0&0\\0&0&1\\0&1&0\end{pmatrix},\qquad
 B_\eps=\begin{pmatrix}0&1&0\\-1&1&0\\1&0&-\eps\end{pmatrix}.
\end{equation}
They satisfy $A_\eps^2=(A_\eps B_\eps)^3=B_\eps^6=I$.

If $\eps=1$, the table gives $T^3=-\mathrm{Id}$ on all of $V$
(equivalently, $B_1^3=-I$).  Lemma~\ref{lem:kernel} immediately implies
$h|\Gamma(6)=h$.  This part has used only hypotheses~(1) and~(2).

Suppose now that $\eps=-1$.  Put
\begin{equation}\label{eq:u}
 u=g+v.
\end{equation}
The table gives
\begin{equation}\label{eq:u-inv}
 u|S=u,\qquad u|T=u.
\end{equation}
Thus $u$ is invariant under the full modular group.  We have not yet
proved $u=0$, and its invariance alone would not suffice.  The next
section uses the Euler condition at $5$ to prove precisely this.

\section{The \texorpdfstring{$m=5$}{m=5} relation kills the defect}

We record the needed prime relation explicitly.  Hypothesis~(3) is
equivalent to
\[
 f|\mathcal T_5=5^{1-k/2}a_5f,\qquad
 \mathcal T_5=
 \begin{pmatrix}5&0\\0&1\end{pmatrix}
 +\sum_{r=0}^4\begin{pmatrix}1&r\\0&5\end{pmatrix}.
\]
Conjugate this equation by $H_{36}$, using \eqref{eq:fricke}, and subtract
the two equations.  The two diagonal terms cancel.  Right-multiplication
by $\operatorname{diag}(5,1)$ then gives
\begin{equation}\label{eq:prime-before}
 f\bigg|\sum_{r=1}^4\bb(r/5)
 =f\bigg|\sum_{c=1}^4\begin{pmatrix}25&0\\-180c&1\end{pmatrix}.
\end{equation}
We also know $f|W_{36}=f$ for
$W_{36}=\begin{pmatrix}1&0\\36&1\end{pmatrix}$, by periodicity and
\eqref{eq:fricke}.  Changing $c$ by a multiple of $5$ on the right-hand
side of \eqref{eq:prime-before} is therefore permitted by left
multiplication with a power of $W_{36}$.  Reindex $c$ by
$36bc\equiv-1\pmod5$, and set
\[
 \gamma_b=\begin{pmatrix}5&-b\\-36c&(1+36bc)/5\end{pmatrix}.
\]
Since
$\gamma_b\bb(b/5)=\begin{pmatrix}5&0\\-36c&1/5\end{pmatrix}$,
we obtain the Conrey--Farmer prime relation
\begin{equation}\label{eq:prime}
 f\bigg|\sum_{b\in\{1,-1,2,-2\}}(1-\gamma_b)\bb(b/5)=0.
\end{equation}
This is the $m=5$ instance of \cite[Corollary~2]{CF}; the derivation above
also checks the normalizations and the permitted left replacements.

Choose explicitly
\[
 \gamma_1=\begin{pmatrix}5&-1\\36&-7\end{pmatrix},\quad
 \gamma_{-1}=\begin{pmatrix}5&1\\-36&-7\end{pmatrix},\quad
 \gamma_2=\begin{pmatrix}5&-2\\-72&29\end{pmatrix},\quad
 \gamma_{-2}=\begin{pmatrix}5&2\\72&29\end{pmatrix}.
\]
For $D=\operatorname{diag}(1,6)$ put $\delta_b=D^{-1}\gamma_bD$.
Conjugating \eqref{eq:prime} gives
\begin{equation}\label{eq:prime-h}
 h\bigg|\sum_b(1-\delta_b)\bb(6b/5)=0.
\end{equation}
The four integral matrices are
\[
 \delta_1=\begin{pmatrix}5&-6\\6&-7\end{pmatrix},\quad
 \delta_{-1}=\begin{pmatrix}5&6\\-6&-7\end{pmatrix},\quad
 \delta_2=\begin{pmatrix}5&-12\\-12&29\end{pmatrix},\quad
 \delta_{-2}=\begin{pmatrix}5&12\\12&29\end{pmatrix}.
\]
The following are exact matrix identities:
\begin{align}
 \delta_{\pm1}&=T^{\pm1}ST^{\pm6}ST^{\mp1},\label{eq:word1}\\
 \delta_{\pm2}&=(ST^{\pm2}ST^{\mp2})^2.\label{eq:word2}
\end{align}
Continue to assume $\eps=-1$.  Since $T^6=1$ on $V$, identity
\eqref{eq:word1} gives $h|\delta_{\pm1}=h$.  For the other two,
\eqref{eq:table}--\eqref{eq:u-inv} give
\[
 h|ST^2ST^{-2}=h-u,\qquad
 h|ST^{-2}ST^2=h+u.
\]
The function $u$ is invariant under both words.  Squaring them proves
\begin{equation}\label{eq:defects}
 h|\delta_2=h-2u,\qquad h|\delta_{-2}=h+2u.
\end{equation}
Thus \eqref{eq:prime-h} reduces to
\[
 2u|\bb(12/5)-2u|\bb(-12/5)=0.
\]
Consequently $u|\bb(24/5)=u$, and $u|T=u$ then implies
$u|\bb(1/5)=u$.  Together with $u|S=u$, this gives
\[
 u|E=u,\qquad E=S\bb(1/5)=\begin{pmatrix}0&-1\\1&1/5\end{pmatrix}.
\]
Here $\det E=1$ and $\operatorname{tr}E=1/5$.  The matrix is elliptic
and has infinite order: a finite-order matrix has algebraic-integer
trace, whereas the rational number $1/5$ is not an algebraic integer.
Lemma~\ref{lem:elliptic} proves $u=0$.

It follows that $v=-g$, so $V=\operatorname{span}\{h,g\}$, with
\[
 h|S=-h,\quad g|S=-g,\quad h|T=g,\quad g|T=g-h.
\]
Again $T^3=-\mathrm{Id}$ on $V$.  Lemma~\ref{lem:kernel} proves
$h|\Gamma(6)=h$ in this sign as well.

\section{Cuspidality and descent to level \texorpdfstring{$36$}{36}}

We have established, for either sign, that $h$ is $\Gamma(6)$-invariant,
that $V=\operatorname{span}\{h|\sigma:\sigma\in\SL_2(\mathbb Z)\}$ is
finite dimensional, and that $T^3=-\mathrm{Id}$ on $V$.
For every $\sigma\in\SL_2(\mathbb Z)$, this gives
\begin{equation}\label{eq:anti-period}
 (h|\sigma)(z+3)=-(h|\sigma)(z).
\end{equation}

We justify cusp holomorphy without adding a Ramanujan-type estimate.
From $a_n=O(n^C)$, uniformly in $x$,
\[
 |f(x+iy)|\ll y^{-C-1}\quad(0<y\le1),
 \qquad |f(x+iy)|\ll e^{-\pi y}\quad(y\ge1).
\]
The first bound follows by comparison of
$\sum n^C e^{-2\pi ny}$ with its defining integral.  If the lower-left
entry of $\sigma$ is nonzero, then on a bounded horizontal strip as
$y\to\infty$, the imaginary part of $\sigma(x+iy)/6$ is comparable to
$y^{-1}$ and the slash factor is $O(y^{-k})$.  Thus $h|\sigma$ grows
at most polynomially at infinity.  If that entry is zero, it decays
exponentially.

Each $h|\sigma$ is holomorphic and $6$-periodic, so it is a holomorphic
function of $q_6=e^{2\pi iz/6}$ on a punctured disk.  Polynomial growth
in $y$ rules out all negative Laurent coefficients: on $|q_6|=r$ the
bound is $O((1+|\log r|)^A)$ for some $A$, and Cauchy's formula bounds
the coefficient of $q_6^{-j}$ by
$O(r^j(1+|\log r|)^A)\to0$ for $j>0$.
The constant coefficient vanishes by \eqref{eq:anti-period}.
Therefore $h\in S_k(\Gamma(6))$.

Finally, for $\gamma=\begin{pmatrix}a&b\\36c&d\end{pmatrix}\in\Gamma_0(36)$,
\[
 D^{-1}\gamma D=\begin{pmatrix}a&6b\\6c&d\end{pmatrix}.
\]
Because $ad\equiv1\pmod6$ and $(\mathbb Z/6\mathbb Z)^\times=\{1,-1\}$,
this belongs to $\Gamma(6)$ or $-\Gamma(6)$.  Conversely every matrix
in $\Gamma(6)\cup-\Gamma(6)$ is obtained in this way.  Thus
\[
 D^{-1}\Gamma_0(36)D=\Gamma(6)\cup-\Gamma(6).
\]
Since $k$ is even and $h=6^{k/2}f|D$, transformation laws and cusp
vanishing transfer under this conjugation.  We obtain
$f\in S_k(\Gamma_0(36))$, completing the proof of
Theorem~\ref{thm:main}.

\section{Verification}

The accompanying file \texttt{verify\_N36.py} uses only exact integer and
rational arithmetic from the Python standard library.  It verifies the
four matrices, their word factorizations, the representation relations,
the two defect formulas, and the conjugation of the prime relation.
Independently of Lemma~\ref{lem:kernel}, it verifies both final
representations on all $72$ elements of
$\SL_2(\mathbb Z/6\mathbb Z)/\{\pm I\}$, checking both generator edges
at every element.  These computations check the algebra; the analytic
vanishing and cusp arguments are proved above.

\begin{thebibliography}{9}
\bibitem{CF}
J. B. Conrey and D. W. Farmer,
\emph{An extension of Hecke's converse theorem},
International Mathematics Research Notices 1995, no.~9, 445--463.
\href{https://arxiv.org/abs/math/9502209}{arXiv:math/9502209}.

\bibitem{CFOS}
J. B. Conrey, D. W. Farmer, B. E. Odgers, and N. C. Snaith,
\emph{A converse theorem for $\Gamma_0(13)$},
Journal of Number Theory 122 (2007), no.~2, 314--323.
\href{https://arxiv.org/abs/math/0601549}{arXiv:math/0601549}.

\bibitem{BCOZ}
B. Bedert, G. Cooper, T. Oliver, and P. Zhang,
\emph{Twisting moduli for $\mathrm{GL}(2)$},
Journal of Number Theory 217 (2020), 142--162.
\href{https://arxiv.org/abs/2003.02557}{arXiv:2003.02557}.
\end{thebibliography}
\end{document}
